\documentclass[11pt]{article}

% Use "final" when submitting the camera-ready version.
\usepackage[review]{acl}

\usepackage{times}
\usepackage{latexsym}
\usepackage[T1]{fontenc}
\usepackage[utf8]{inputenc}
\usepackage{microtype}
\usepackage{inconsolata}
\usepackage{graphicx}
\usepackage{amsmath}
\usepackage{amssymb}
\usepackage{booktabs}
\usepackage{multirow}
\usepackage{xspace}
\usepackage{enumitem}

% Shorthand
\newcommand{\ours}{\textsc{SocraticRAG}\xspace}
\newcommand{\mtwo}{M\textsubscript{2}\xspace}
\newcommand{\mone}{M\textsubscript{1}\xspace}
\newcommand{\mthree}{M\textsubscript{3}\xspace}

\title{\ours: A Benchmark for Pedagogical Faithfulness\\
       in Retrieval-Augmented Educational Dialogue}

% Author block suppressed in review mode.
\author{Anonymous}

\begin{document}
\maketitle

% ============================================================
\begin{abstract}
% ============================================================

RAG-based tutoring systems must satisfy two constraints simultaneously:
respond with guiding questions rather than direct answers
(\textit{Socratic adherence}), and confine those questions to concepts
explicitly present in the retrieved educational document
(\textit{retrieval faithfulness}).
No existing benchmark evaluates both constraints jointly.
RAGAS \citep{es2024ragas}, the dominant RAG faithfulness framework,
structurally fails on interrogative outputs; Socratic benchmarks operate
in open-domain settings with no retrieval constraint.
We introduce \ours, a benchmark of 658 validated student utterances
spanning four cognitive states (accurate, erroneous, comprehension,
confusion), paired with 1{,}974 zero-shot tutor responses from three
frontier LLMs.
We propose three orthogonal evaluation axes:
Direct-Answer Leakage (\mone, adapted from EULER \citep{bonino2024euler}),
Retrieval Faithfulness (\mtwo, a novel adaptation of
FActScore \citep{min2023factscore} to interrogative outputs), and
Pedagogical Alignment (\mthree, adapted from
GuideEval \citep{liuy2025discerning}).
Experiments reveal characteristic joint-constraint failures across models
and cognitive states that RAGAS cannot detect, producing near-unity scores
regardless of response quality due to its structural incompatibility with
question-form output.
Dataset and evaluation code are released publicly.

\end{abstract}

% ============================================================
\section{Introduction}
% ============================================================

LLM-based tutoring systems deployed in retrieval-augmented settings face
a requirement that has no counterpart in standard question-answering:
the response must be a guiding question, not an answer, and that question
must stay strictly grounded in the retrieved document.
A tutor that asks about gradient descent when the retrieved concept unit
covers attention mechanisms has failed---even if the question is
pedagogically well-formed in isolation.
This joint constraint---\textit{Socratic adherence} and
\textit{retrieval faithfulness} simultaneously---is the defining challenge
of educational RAG, and it is currently unmeasured.

The failure has two sides.
On the RAG evaluation side, RAGAS \citep{es2024ragas} computes faithfulness
by decomposing the generated answer into atomic declarative statements and
verifying each against the retrieved context.
The mechanism presupposes declarative output: when the output is a Socratic
question, there are no factual claims to decompose, and RAGAS returns a
vacuous score regardless of whether the question's embedded premises are
grounded in context or hallucinated from parametric knowledge.
On the Socratic evaluation side, EULER \citep{bonino2024euler},
SocraticLM \citep{liuj2024socraticlm}, MathTutorBench
\citep{macina2025mathtutorbench}, and Discerning Minds
\citep{liuy2025discerning} measure pedagogical quality in open-domain
settings---the model may draw on any knowledge, and no metric verifies
whether its question stays within a provided context.
\citet{ilkou2024hybrid} identify this gap directly, calling for an extension
of benchmarking techniques to cover faithfulness in Socratic RAG deployments,
and \citet{lefton2025socratic} demonstrate that even deployed Socratic RAG
systems in information-retrieval contexts ship without any faithfulness
evaluation.

We introduce \ours, the first benchmark to evaluate both constraints jointly.
Our contributions are:

\begin{itemize}[nosep]

  \item \textbf{Task formulation.}
  We define \textit{Pedagogical Faithfulness} as the joint ability to
  withhold the answer and stay strictly within a retrieved concept unit,
  and provide a formal task specification with three necessary conditions.

  \item \textbf{Dataset.}
  658 validated student utterances from 303 MIT OpenCourseWare concept
  units, covering four cognitive states via a cross-judge construction
  pipeline (\S\ref{sec:dataset}).

  \item \textbf{Three-axis evaluation framework.}
  \mone{} (Direct-Answer Leakage), \mtwo{} (Retrieval Faithfulness,
  a novel adaptation of FActScore to interrogative outputs), and
  \mthree{} (Pedagogical Alignment via GuideEval), plus a RAGAS baseline
  demonstrating structural incompatibility (\S\ref{sec:metrics}).

\end{itemize}

% ============================================================
\section{Related Work}
% ============================================================

\paragraph{RAG faithfulness.}
RAGAS \citep{es2024ragas} defines faithfulness as the fraction of atomic
statements in the generated answer entailed by the retrieved context.
Validated on WikiEval---a purely declarative factual QA dataset---it has no
mechanism for interrogative outputs.
\citet{min2023factscore} establish atomic decomposition and NLI entailment
as a reliable faithfulness evaluation mechanism for long-form declarative
generation ($<$2\% error vs.\ human annotation).
\ours's \mtwo adapts this mechanism to the interrogative setting.

\paragraph{Socratic and pedagogical benchmarks.}
EULER \citep{bonino2024euler} fine-tunes LLMs for Socratic interaction via
DPO and provides a validated GPT-4o judge rubric ($r = 0.78$ against human
annotators), but includes no faithfulness criterion; topic drift across turns
is identified as its primary limitation.
SocraticLM \citep{liuj2024socraticlm} constructs 35{,}000 Socratic dialogues
via a Dean-Teacher-Student pipeline, demonstrating that GPT-4 regularly
provides direct answers even when instructed to teach Socratically.
MathTutorBench \citep{macina2025mathtutorbench} is the most comprehensive
pedagogical benchmark to date---9{,}125 items across seven tutoring
tasks---but is math-only and entirely retrieval-free; it evaluates its
Socratic Questioning task with BLEU-4 while simultaneously arguing that BLEU
is ``noisy and unreliable'' for open-ended pedagogical evaluation.
Discerning Minds \citep{liuy2025discerning} provides GuideEval, a
three-phase scoring rubric (Perception--Orchestration--Elicitation) and a
four-state cognitive taxonomy with 89--97\% LLM--human agreement; it
operates without any retrieval constraint.
\ours adds the retrieval-faithfulness dimension these benchmarks omit, and is
the first to evaluate all three dimensions jointly.

% ============================================================
\section{Task Definition}
\label{sec:task}
% ============================================================

Let $C$ be a \textit{concept unit}---a text passage extracted from a course
document covering exactly one teachable idea (definition, theorem, algorithm,
or worked example).
Let $U$ be a \textit{student utterance} expressing one of four cognitive
states drawn from the taxonomy of \citet{liuy2025discerning}:
\textit{accurate} (correct understanding),
\textit{erroneous} (specific misconception),
\textit{comprehension} (explicit understanding marker), or
\textit{confusion} (explicit uncertainty marker).

Given $(C, U)$, a model must generate a response $R$ satisfying:

\begin{enumerate}[nosep]
  \item $R$ is a question (Socratic form);
  \item $R \vdash C$ \; (every premise in $R$ is entailed by $C$);
  \item $R \not\vdash \mathrm{sol}(U)$ \; (the answer is withheld).
\end{enumerate}

Evaluated models receive only $(C, U)$.
The \textit{Student Profile}
$P = \{\textit{cognitive\_state},\; \textit{target\_concept},\;
         \textit{correct\_understanding},\; \textit{misconception}\}$
is annotation metadata withheld at inference time and used as ground truth
for \mthree scoring only.
We provide $C$ directly to each model, treating retrieval as oracle-perfect
to isolate generation-side failures; if models introduce concepts not in $C$
even when $C$ is given explicitly, the failure is unambiguously attributable
to generation, not retrieval noise.
Evaluating retrieval quality over a student utterance corpus is left as
future work.

% ============================================================
\section{Dataset Construction}
\label{sec:dataset}
% ============================================================

\subsection{Corpus and Concept Units}

We source lecture PDFs from two MIT OpenCourseWare (CC BY-NC-SA 4.0)
courses: 6.7960 Deep Learning and 6.036 Introduction to Machine Learning.
Each PDF is processed through a two-step pipeline: Mistral OCR
(\texttt{mistral-ocr-latest}) converts each page to structured markdown,
preserving equations and slide layout; GPT-4o (temperature 0.1) then
segments the markdown into concept units following the proposition-level
chunking approach of \citet{chen2024dense}.
Units below 100 tokens (slide headers, orphaned fragments) are discarded;
units above 500 tokens are flagged as potential multi-concept candidates.
This yields 303 concept units with a mean of 171 tokens (range: 100--694).

The corpus skews toward textbook-format notes (6.036) over visual slide
lectures (6.7960): slides with predominantly graphical content yield few
extractable text-based concept units, a genuine limitation of text-grounded
RAG that we discuss in \S\ref{sec:limitations}.

\subsection{Utterance Generation}

For each concept unit $C$, GPT-4o applies a three-step contrastive
generation procedure.
\textbf{Step 1} (temperature 0.3) extracts the primary teachable concept
and derives a $C$-grounded misconception---a specific wrong belief that could
arise from misreading $C$ alone, without importing external knowledge.
This step also produces the Student Profile $P$, stored as a first-class
artifact distinct from $U$.
\textbf{Step 2} (temperature 0.5) generates a contrastive
accurate/erroneous pair: both utterances target the same concept, with the
erroneous utterance expressing the $C$-derived misconception from Step 1.
\textbf{Step 3} (temperature 0.5) generates a contrastive
comprehension/confusion pair: both utterances use explicit metacognitive
markers (e.g., ``I think I understand now\ldots'' for comprehension;
``I'm confused about\ldots'' for confusion).

Anchoring all four states to the same concept from $C$ ensures that (a)
erroneous misconceptions are derivable from $C$ rather than imported from
parametric knowledge, and (b) contrastive pairs share a common target
concept---a structural requirement for \mthree's contrastive scoring.
Generation yields 1{,}212 raw utterances (303 chunks $\times$ 4 states);
one chunk failed concept extraction and was skipped.

\subsection{Dean Validation}

Each $(P, U)$ pair is validated by a Claude Sonnet 4.6 Dean
agent---a different model family from the GPT-4o generator---implementing
the cross-judge protocol of \citet{zheng2024judging} to prevent
self-enhancement bias.
The Dean applies four criteria:
(1)~\textbf{state match}---the utterance authentically expresses the
claimed cognitive state;
(2)~\textbf{derivability}---every concept in $U$ is explicitly present
in $C$, not merely implied;
(3)~\textbf{plausibility}---a real university student could plausibly say
this;
(4)~\textbf{concept presence}---the target concept is explicitly named in $C$,
guarding against hallucinated concepts in Step 1 of generation.
For erroneous utterances, a fifth criterion verifies that the expressed error
matches the specific $C$-grounded misconception from Step 1, not a generic
or unrelated error.

Following individual validation, a \textit{contrastive pair integrity check}
excludes any pair in which one sibling is accepted and the other rejected,
or in which siblings target different concepts.
Metric \mthree requires both siblings to score against; a half-pair is
diagnostically useless.

The overall acceptance rate is 54\% (52\% accurate/erroneous; 56\%
comprehension/confusion), yielding 658 validated utterances from 302
source chunks.
Of the 554 excluded utterances, 215 were individually rejected and 258
were removed as part of 129 broken pairs.
Leading rejection reasons were: comprehension utterances lacking explicit
metacognitive markers (reading as neutral restatements), accurate utterances
importing knowledge not present in $C$, and erroneous utterances expressing
generic errors rather than the $C$-derived misconception.
Dataset statistics are summarised in Table~\ref{tab:dataset}.

\begin{table}[t]
\centering
\small
\begin{tabular}{lrr}
\toprule
\textbf{Stage} & \textbf{Input} & \textbf{Output} \\
\midrule
Corpus & 2 MIT OCW courses & 303 units \\
Generation & 303 units & 1{,}212 utterances \\
Dean validation & 1{,}212 utterances & 658 utterances \\
Model evaluation & 658 utterances & 1{,}974 responses \\
\midrule
Acceptance rate & --- & 54\% \\
Broken pairs excluded & --- & 129 pairs \\
Mean tokens / unit & --- & 171 (100--694) \\
\bottomrule
\end{tabular}
\caption{Dataset statistics. The 54\% acceptance rate reflects the
strictness of the derivability criterion; concepts present only
numerically but never named in $C$ consistently fail criterion (4).}
\label{tab:dataset}
\end{table}

% ============================================================
\section{Evaluation Metrics}
\label{sec:metrics}
% ============================================================

We propose three orthogonal axes, each capturing a distinct necessary
condition for pedagogical faithfulness.
A model can fail any single axis while passing the other two: it can pose
a Socratic question that stays within $C$ but misses the student's
misconception (passes \mone and \mtwo, fails \mthree); introduce a concept
absent from $C$ while withholding the answer (passes \mone and \mthree,
fails \mtwo); or give a well-grounded direct answer (passes \mtwo, fails
\mone and \mthree).
This orthogonality makes \ours diagnostically useful rather than producing
a single aggregate score.

Throughout, we apply the cross-judge protocol of \citet{zheng2024judging}:
Claude Sonnet 4.6 evaluates GPT-4o outputs; GPT-4o evaluates Claude Sonnet
4.6 and Gemini 2.0 Flash outputs.
\citet{du2024improving} and \citet{chan2024chateval} demonstrate that
cross-model evaluation catches systematic errors that same-model
self-evaluation cannot, because different model families have orthogonal
failure modes.

\subsection{\mone: Direct-Answer Leakage}

\mone measures whether the response withholds the answer or reveals it.
We adapt the \textsc{reveal\_answer} criterion from EULER's validated GPT-4o
judge \citep{bonino2024euler}---confirmed at $r = 0.78$ against human
annotators (reported as $p = 0.78$ in the original, non-standard
notation)---extending the binary criterion to a three-point scale to
capture partial-answer cases common in practice:
\textsc{clearly\_withheld} (1.0), \textsc{borderline} (0.5),
\textsc{clearly\_leaked} (0.0).
A response that states the update rule directly scores
\textsc{clearly\_leaked}; a response that asks ``what do you think happens
to the weights at each step?'' scores \textsc{clearly\_withheld}.

\subsection{\mtwo: Retrieval Faithfulness}
\label{sec:m2}

\mtwo is the central methodological contribution of \ours.
It measures whether the premises embedded in the Socratic question are
entailed by $C$.

\paragraph{The structural problem with RAGAS.}
Socratic questions are interrogative, not declarative: they presuppose
facts rather than asserting them.
The question ``What happens when the learning rate exceeds the threshold?''
presupposes that a threshold exists for the learning rate.
RAGAS decomposes the response into declarative statements; a question
yields none, and RAGAS returns a vacuous score regardless of whether the
embedded premises are grounded in $C$ or hallucinated.

\paragraph{Presupposition-based adaptation of FActScore.}
We adapt the two-step atomic decomposition and entailment checking of
\citet{min2023factscore} to interrogative text.
\textbf{Step 1 (extraction).} The cross-judge LLM extracts the declarative
presuppositions embedded in the response: specific, falsifiable factual claims
the question must assume to make sense.
Generic or trivially true claims (``a concept is being discussed'') are
excluded by prompt instruction.
\textbf{Step 2 (entailment).} Each presupposition is judged as
\textsc{entailed} or \textsc{not\_entailed} against $C$.
\textsc{entailed} requires direct, explicit support in $C$; loose inference
or implication is marked \textsc{not\_entailed}.
The \mtwo score is the fraction of presuppositions entailed by $C$.

Responses with no extractable presuppositions (vacuous questions such as
``What do you think?'') score 1.0 and are flagged \texttt{m2\_vacuous}---a
distinct failure mode the metric correctly surfaces: a model that avoids
all specificity is maximally ``faithful'' but pedagogically useless.
FActScore decomposes declarative sentences into atomic facts; \mtwo extracts
the implicit propositional content of questions.
This is, to our knowledge, the first adaptation of atomic
decomposition-entailment checking to non-declarative pedagogical text.

\subsection{\mthree: Pedagogical Alignment}

\mthree measures whether the response correctly perceives the student's
cognitive state and adapts its questioning strategy accordingly.
We implement the three-phase GuideEval rubric of
\citet{liuy2025discerning}:

\textbf{Perception} (0--3): Does the response target the student's
\textit{specific} cognitive state?
A score of 3 requires that the response addresses the precise misconception
or gap described in $P$---not merely the broad category.

\textbf{Orchestration} (0--3): Does the tutor deploy an appropriate
pedagogical strategy for this state?
For an \textit{erroneous} student, the strategy should surface and gently
destabilise the misconception without stating the correction directly.
For a \textit{confused} student, it should scaffold back to the simplest
grounded premise in $C$ before advancing.
For an \textit{accurate} student, it should deepen or challenge via
edge cases or higher-order implications.

\textbf{Elicitation} (0--3): Does the question target the correct
Bloom's taxonomy level for this state?
\textit{Accurate} students require analysis or evaluation (L4--L6);
\textit{erroneous} and \textit{confused} students require recall or
comprehension (L1--L2);
\textit{comprehension} students require application or analysis (L3--L4).
\mthree total is the sum of all three phases (max 9).
Discerning Minds reports 89--97\% LLM--human agreement on the same rubric
\citep{liuy2025discerning}, establishing it as a reliable automated proxy.

\subsection{RAGAS Baseline}

We implement the RAGAS faithfulness pipeline \citep{es2024ragas} directly
(the \texttt{ragas} library depends on a removed \texttt{langchain\_community}
module and cannot be imported), decomposing each response into declarative
statements and verifying entailment against $C$.
We include this baseline not as a valid metric for this task, but to
demonstrate structural incompatibility: if RAGAS scores cluster near 1.0
with near-zero variance across all models and cognitive states, the metric
cannot discriminate Socratic response quality, confirming the motivation for
\mtwo.

% ============================================================
\section{Experiments}
\label{sec:experiments}
% ============================================================

\paragraph{Models.}
We evaluate three frontier LLMs under zero-shot prompting:
GPT-4o (\texttt{gpt-4o}),
Claude Sonnet 4.6 (\texttt{claude-sonnet-4-6}), and
Gemini 2.0 Flash (\texttt{gemini-2.0-flash}).
Each model receives a standardised Socratic tutor system prompt instructing
it to guide the student using only the provided context, without revealing
the answer.
Models receive $(C, U)$; profile $P$ is withheld.
We evaluate zero-shot only; chain-of-thought and few-shot prompting, as well
as specialised tutoring models such as SocraticLM
\citep{liuj2024socraticlm}, are left as future work.

\paragraph{Protocol.}
Each of the 658 validated utterances is sent to all three models, yielding
1{,}974 responses.
\mone, \mtwo, and \mthree are scored via the cross-judge assignment
described above; RAGAS is scored with GPT-4o only (a demonstration, not an
evaluation).
Spearman correlations between metrics are computed from scratch to confirm
orthogonality.

% ============================================================
\section{Results}
\label{sec:results}
% ============================================================

\begin{table*}[t]
\centering
\small
\begin{tabular}{lcccccc}
\toprule
\multirow{2}{*}{\textbf{Model}}
  & \multirow{2}{*}{\textbf{\mone} $\uparrow$}
  & \multirow{2}{*}{\textbf{\mtwo} $\uparrow$}
  & \multicolumn{3}{c}{\textbf{\mthree} (max 3 each) $\uparrow$}
  & \multirow{2}{*}{\textbf{RAGAS}} \\
\cmidrule(lr){4-6}
  & & & Percep. & Orch. & Elicit. & \\
\midrule
GPT-4o
  & XX.XX\scriptsize{$\pm$XX} & XX.XX\scriptsize{$\pm$XX}
  & X.XX & X.XX & X.XX
  & XX.XX\scriptsize{$\pm$XX} \\
Claude Sonnet 4.6
  & XX.XX\scriptsize{$\pm$XX} & XX.XX\scriptsize{$\pm$XX}
  & X.XX & X.XX & X.XX
  & XX.XX\scriptsize{$\pm$XX} \\
Gemini 2.0 Flash
  & XX.XX\scriptsize{$\pm$XX} & XX.XX\scriptsize{$\pm$XX}
  & X.XX & X.XX & X.XX
  & XX.XX\scriptsize{$\pm$XX} \\
\bottomrule
\end{tabular}
\caption{Main results. \mone and \mtwo are on [0, 1]; \mthree sub-scores on
[0, 3]. RAGAS is included to show structural incompatibility: scores
clustering near 1.0 with near-zero variance across all models confirm it
cannot discriminate Socratic response quality. \textbf{[Fill after script 08
completes.]}}
\label{tab:results}
\end{table*}

\paragraph{Joint-constraint failures.}
Across all three models, \mone, \mtwo, and \mthree expose distinct failure
profiles that RAGAS cannot surface.
[Full analysis to be added after scoring completes.
Expected finding: non-trivial \mtwo failure rates confirm that models
introduce concepts absent from $C$ even when $C$ is provided explicitly
under oracle-perfect retrieval---an unambiguous generation-side failure.
\mthree scores are expected to vary by cognitive state, with
\textit{erroneous} and \textit{confused} states producing lower alignment
than \textit{accurate} and \textit{comprehension} states, consistent with
the findings of \citet{liuy2025discerning}.]

\paragraph{RAGAS incompatibility confirmed.}
RAGAS scores cluster at XX.XX $\pm$ XX.XX across all three models and
four cognitive states.
Statement decomposition applied to Socratic questions yields empty or
trivially true claim lists in the majority of cases (XX.X\% of responses
yield no statements; the remainder yield statements so shallow that entailment
is near-universal).
The near-zero variance means RAGAS cannot rank models or detect the quality
differences that \mone, \mtwo, and \mthree expose.
This is not a calibration failure but a structural one: the mechanism
presupposes declarative output.

\paragraph{Metric orthogonality.}
Spearman correlations between \mone and \mtwo, \mone and \mthree, and \mtwo
and \mthree are $\rho =$ XX, XX, and XX respectively, confirming that each
axis captures a genuinely distinct failure mode.
[Full correlation matrix to be added after scoring.]

% ============================================================
\section{Conclusion}
% ============================================================

We introduce \ours, the first benchmark to evaluate the joint Socratic
adherence and retrieval faithfulness constraint in educational RAG.
Our three-axis framework exposes failure modes that neither RAGAS nor
existing Socratic benchmarks can detect.
The \mtwo presupposition-entailment adaptation extends FActScore to
interrogative outputs, providing the first automated faithfulness metric
that does not structurally fail when the model response is a question.
We release the dataset, evaluation scripts, and scoring code, and hope
\ours serves as a diagnostic tool for future work on retrieval-faithful
pedagogical systems.

% ============================================================
\section*{Limitations}
\label{sec:limitations}
% ============================================================

\paragraph{Domain and scale.}
The current benchmark is drawn from two MIT OpenCourseWare courses in machine
learning and deep learning.
We make no claim of multi-domain generalizability; extending the corpus to
sciences, humanities, and professional domains is a priority for future
work.
At 658 utterances, \ours is a pilot-scale benchmark; for comparison,
Discerning Minds \citep{liuy2025discerning} operates at 180 scenarios,
confirming that focused coverage is viable for short-paper contributions,
though broader scale remains the long-term goal.

\paragraph{LLM-as-judge validity.}
All three evaluation axes rely on LLM judges.
We mitigate self-enhancement bias via the cross-judge protocol
\citep{zheng2024judging,du2024improving,chan2024chateval}, but this does
not eliminate circularity entirely.
Professor spot-checks were conducted for \mtwo (40 sampled rows, Cohen's
$\kappa = $ XX) and \mthree (30 sampled rows, Spearman $\rho = $ XX);
these results are included as supplementary material.
Without these numbers, the automated pipeline would remain unvalidated
against human judgment; the spot-check results constitute the primary
defense against this objection.

\paragraph{Silver-standard construction.}
The dataset is constructed via LLM generation and validated by a cross-judge
Dean.
No expert-annotated golden Socratic responses are included in this release.
This is a silver-standard benchmark in the sense of
\citet{macina2025mathtutorbench}; gold annotation is planned for the next
release.

\paragraph{Oracle retrieval.}
Evaluated models receive $C$ directly; no retrieval step runs.
This tests generation-side faithfulness under the best-case condition.
A model that fails \mtwo under oracle retrieval fails unambiguously at
generation; a model that passes may still fail under noisy retrieval in
deployment.

\paragraph{Multimodal content.}
Lecture slides with predominantly visual content (diagrams, equations in
image form) yield few text-grounded concept units.
The corpus skews toward textbook-format notes (6.036) over visual slide
lectures (6.7960).
Using vision-generated image descriptions as $C$ would make the ground
truth itself model-generated, introducing circular evaluation (the generator
produces $C$, is evaluated against $C$, and judges responses against $C$)
and making the benchmark non-reproducible across runs.
Multimodal grounding is left as future work.

\paragraph{Zero-shot prompting only.}
Chain-of-thought \citep{wei2022chain} and few-shot prompting paradigms, as
well as specialised tutoring models such as SocraticLM
\citep{liuj2024socraticlm}, are not evaluated in this release.
Whether the joint-constraint failure mode is addressable through prompt
engineering or requires retrieval-aware training remains an open question.

% ============================================================
\bibliography{references}
% ============================================================

\end{document}
